\documentclass[10pt,a4paper]{amsart}
\usepackage[margin=19mm]{geometry}
\usepackage{amsmath,amssymb,mathtools}
\usepackage[scr=rsfs,cal=euler]{mathalfa}
\usepackage{xfrac}
\usepackage[unicode,hidelinks,bookmarks=false]{hyperref}
\newcommand{\ie}{i.e.\ }
\newcommand{\cf}{cf.\ }
\newcommand{\resp}{respectively\ }
\newcommand{\Subsection}{\subsection}
\providecommand{\nobreakdash}{\nobreak\hbox{-}}
\setlength{\emergencystretch}{2em}
\multlinegap0pt
\usepackage{amssymb}
\newtheorem{lemma}{Lemma}
\title{Strassen's local law of the iterated logarithm for the generalized fractional Brownian motion}
\author{Ran Wang, Yimin Xiao}
\date{Source: arXiv:2411.15681v2 (CC BY 4.0)}
\begin{document}
\maketitle

\noindent\emph{Source and licence.} Strassen's local law of the iterated logarithm for the generalized fractional Brownian motion. Version arXiv:2411.15681v2 (CC BY 4.0), distributed by arXiv under the Creative Commons Attribution 4.0 International licence, http://creativecommons.org/licenses/by/4.0/, as recorded in the arXiv licence metadata for this exact version and shown on the abstract page https://arxiv.org/abs/2411.15681v2. Full licence notice: \texttt{licenses/arxiv-2411.15681-CC-BY-4.0.txt}.\par\smallskip

\noindent\textit{Editorial scope.} Counted unit: the article's lemma const (source0.tex lines 1100--1111). The article's complete original proof of it is delivered with it (source0.tex lines 1112--1167, one \texttt{proof} environment of 56 lines). No mathematics is written, repaired or re-derived: the statement and the proof are the article's own lines, in the article's own notation, and no retained line is substituted or re-flowed.  Retained uncounted as the source-local context the proof needs: lines 1069--1071; lines 1073--1075; lines 1077--1079.  Nothing was omitted from any retained span, so the package carries no omission record.  No retained line contains a citation, so the wrapper carries no bibliography and no bibliography entry.  No retained line contains a figure environment, an image-inclusion command or drawing code, so no image is delivered and the package declares no asset dependency.  Editorial formatting only, in addition: an AMS \texttt{amsart} preamble that restates the article's own declarations and macros used by the retained lines, namely the environments \texttt{lemma} on separate counters, together with the standard packages those lines need.  The retained environments print in this document as Lemma 1 in order of appearance, whereas the article prints the same items with the numbering of its own full document.  The complete native source of the article as distributed in the arXiv e-print for this exact version is bundled unchanged as \texttt{sources/169012/source0.tex}.
 \begin{align}\label{eq cov 1}  {
R_H(t, s)=\mathbb E\left[B^H(t) B^H(s) \right]=\frac12\left(|t|^{2H}+|s|^{2H}-|t-s|^{2H}\right), \ \ \ s, t>0. }
 \end{align} {

 \begin{align}\label{eq cov 2}  {
 R_H(t, s)=\int_0^{t\wedge s}K_H(t, u)K_H(s, u)du. }
 \end{align}  {

 \begin{align}\label{eq KH1} {
 K_H(t, u)=\mathbf 1_{[0,t]}(u).  }
 \end{align} {

\begin{lemma}\label{lem const}
\begin{itemize}
\item[(a)]  { It holds that }
\begin{align*}  {
\sup_{\xi \in \mathscr S}\xi(1) =\sup_{\xi \in \mathscr S}\sup_{0\le t\le 1}|\xi(t)|=1.  }
\end{align*}
\item[(b)]   {For $0<\delta<1$, we have }
\begin{align*}  {
\sup_{\xi \in \mathscr S}\sup_{0\le t\le 1-\delta}|\xi(t+\delta)-\xi(t)|   = \sup_{\xi \in \mathscr S}\sup_{0\le t\le 1-\delta}\sup_{0\le s\le \delta}|\xi(t+s)-\xi(t)| \in \left[\delta^H, \sqrt{2}\delta^H\right]. }
\end{align*}
 \end{itemize}
\end{lemma}

 \begin{proof}
\noindent  {  (a)  For any $\xi \in \mathscr S$, by the Cauchy-Schwarz inequality and \eqref{eq cov 2}, we have 
\begin{align*}
\xi(1)=\int_0^1 K_H(1, u) \dot \xi(u)du\le &\,   \sqrt{ \int_0^1 K_H(1, u)^2  du \cdot \int_0^1  |\dot \xi(u)|^2du }\\
 \le & \,  \sqrt{ \int_0^1 K_H(1, u)^2  du} = 1. 
\end{align*}
Moreover, the above equalities holds if we choose $\dot \xi(u)=K_H(1, u)$ on   $[0,1]$.  Thus,  $$\sup_{\xi \in \mathscr S}\xi(1)=1.$$
}

 {
Similarly,   for any $\xi \in \mathscr S$, by the Cauchy-Schwarz inequality and \eqref{eq cov 2}, we have  that for any $t\in [0, 1]$,
\begin{align*}
|\xi(t)|=\left|\int_0^t K_H(t, u) \dot \xi(u)du\right|
\le &\,   \sqrt{ \int_0^t K_H(t, u)^2  du \cdot \int_0^t  |\dot \xi(u)|^2du }\\
\le &\,  \sqrt{ \int_0^t K_H(t, u)^2  du}\le 1. 
\end{align*}
If we choose $\dot \xi(u)=K_H(1, u)$ on $[0,1]$, then we have  $$\sup_{\xi \in \mathscr S}\sup_{0\le t\le 1}|\xi(t)|=1.$$
}
\noindent
  (b)  {Given $0<\delta<1$ and $\xi\in\mathscr S$,  by the Cauchy-Schwarz inequality and \eqref{eq cov 2}, we have that for any $t\in [0, 1-\delta]$, 
 \begin{align*}
 |\xi(t+\delta)-\xi(t)|= &\, \left|\int_0^{t+\delta} K_H(t+\delta, u) \dot \xi(u)du-\int_0^{t } K_H(t, u) \dot \xi(u)du \right|\\
 = &\, \left|\int_0^{t}\left( K_H(t+\delta, u)-K_H(t, u)\right)  \dot \xi(u)du \right| +\left|\int_t^{t+\delta} K_H(t+\delta, u) \dot \xi(u)du  \right|\\
 \le &\,  \left(\int_0^{t}\left( K_H(t+\delta, u)-K_H(t, u)\right)^2 du\right)^{\frac12}+  \left(\int_{t}^{t+\delta} K_H(t+\delta, u) ^2 du\right)^{\frac12}.
\end{align*}
}
 { Moreover,  the equality in the last step is achieved when
 \begin{equation*} 
  \dot\xi(u)= \left\{\begin{array}{ll}
K_H(t+\delta, u)-K_H(t, u),  \ \quad &\hbox{ when }\  u\in [0, t],\\
 K_H(t+\delta, u),  \ \quad &\hbox{ when } \ u\in (t, t+\delta].
\end{array}
\right.
 \end{equation*}
 }
 {By \eqref{eq cov 1} and \eqref{eq cov 2}, we know 
 $$
 \int_0^{t}\left( K_H(t+\delta, u)-K_H(t, u)\right)^2 du +  \int_{t}^{t+\delta} K_H(t+\delta, u) ^2 du=\delta^{2H}.
 $$
 Consequently, using the elementary inequality
  $$ \sqrt{a+b}\le \sqrt{a}+\sqrt{b}\le \sqrt{2(a+b)}, \ \  a, b\ge0, $$
 we have
 $$
\left(\int_0^{t}\left( K_H(t+\delta, u)-K_H(t, u)\right)^2 du\right)^{\frac12}+  \left(\int_{t}^{t+\delta} K_H(t+\delta, u) ^2 du\right)^{\frac12} \in \left[\delta^H, \sqrt{2}\delta^H\right].
 $$
 A parallel argument yields
  $$
   \sup_{\xi \in \mathscr S}\sup_{0\le t\le 1-\delta}\sup_{0\le s\le \delta}|\xi(t+s)-\xi(t)| \in  \left[\delta^H, \sqrt{2}\delta^H\right].
  $$
  In particular,   when $H=\frac12$,   substituting \eqref{eq KH1} into the above estimates leads to the following  exact value:
  \begin{align*}
\sup_{\xi \in \mathscr S}\sup_{0\le t\le 1-\delta}|\xi(t+\delta)-\xi(t)|   = \sup_{\xi \in \mathscr S}\sup_{0\le t\le 1-\delta}\sup_{0\le s\le \delta}|\xi(t+s)-\xi(t)| = \delta^{\frac12}.
\end{align*}
 }
  {The proof is complete. }
 \end{proof}

\end{document}
